\documentclass[12pt, a4paper, twoside]{book}
\usepackage[a4paper, inner=3.0cm, outer=2.5cm, top=2.8cm, bottom=2.8cm, includeheadfoot, headheight=24pt]{geometry}
\usepackage{amsmath, amsthm}
\usepackage[charter, cal=cmcal]{mathdesign}
\usepackage[scaled=0.92]{helvet}
\usepackage{mathtools}
\usepackage{physics}
\usepackage{graphicx}
\usepackage{xcolor}
\usepackage{tikz}
\usetikzlibrary{arrows.meta, calc, positioning, decorations.pathmorphing}
\usepackage{pgfplots}
\pgfplotsset{compat=1.18}
\usepackage{tcolorbox}
\tcbuselibrary{skins, breakable, theorems}
\usepackage{enumitem}
\usepackage{booktabs}
\usepackage{fancyhdr}
\usepackage{titlesec}
\usepackage{hyperref}
\usepackage{float}
\usepackage{caption}
\usepackage{epigraph}
\usepackage{lettrine}
\usepackage{setspace}
\usepackage{parskip}
\usepackage{siunitx}
\usepackage{bm}
\usepackage{array}

\definecolor{deep}{HTML}{1a1a2e}
\definecolor{accent}{HTML}{e94560}
\definecolor{teal}{HTML}{2ec4b6}
\definecolor{gold}{HTML}{f5a623}
\definecolor{midblue}{HTML}{0f3460}
\definecolor{purple}{HTML}{533483}
\definecolor{orange}{HTML}{ff6b35}
\definecolor{green}{HTML}{2d6a4f}
\definecolor{lightbg}{HTML}{f7f9fc}
\definecolor{warmbg}{HTML}{fffbf2}

\setlength{\headheight}{24pt}
\hypersetup{colorlinks=true,linkcolor=midblue,urlcolor=accent,citecolor=teal,
  pdftitle={Discovery Physics J6 M4}}

\pagestyle{fancy}\fancyhf{}
\fancyhead[LE]{\small\textcolor{deep}{\textit{Discovery Physics}}}
\fancyhead[RO]{\small\textcolor{deep}{\textit{\leftmark}}}
\fancyfoot[C]{\small\textcolor{deep}{\thepage}}
\renewcommand{\headrulewidth}{0.4pt}
\renewcommand{\headrule}{\hbox to\headwidth{%
  \color{teal}\leaders\hrule height\headrulewidth\hfill}}

\titleformat{\chapter}[display]
  {\normalfont\huge\bfseries\color{deep}}
  {\textcolor{accent}{\chaptertitlename\ \thechapter}}{12pt}{\Huge}
  [\vspace{4pt}\textcolor{teal}{\titlerule[1.5pt]}]
\titleformat{\section}{\normalfont\Large\bfseries\color{midblue}}
  {\textcolor{accent}{\thesection}}{1em}{}
\titleformat{\subsection}{\normalfont\large\bfseries\color{deep}}
  {\textcolor{teal}{\thesubsection}}{1em}{}
\titleformat{\subsubsection}{\normalfont\normalsize\bfseries\color{midblue}}
  {\thesubsubsection}{1em}{}

\newtcolorbox{puzzlebox}[1][]{enhanced,breakable,colback=warmbg,colframe=gold,
  fonttitle=\bfseries\large\color{deep},title={The Puzzle},
  attach boxed title to top left={yshift=-2mm,xshift=4mm},
  boxed title style={colback=gold},#1}
\newtcolorbox{explobox}[1][]{enhanced,breakable,colback=lightbg,colframe=midblue,
  fonttitle=\bfseries\color{white},coltitle=white,title={Exploration},
  attach boxed title to top left={yshift=-2mm,xshift=4mm},
  boxed title style={colback=midblue},#1}
\newtcolorbox{discoverybox}[1][]{enhanced,breakable,colback=white,colframe=accent,
  fonttitle=\bfseries\large\color{white},coltitle=white,title={Discovery},
  attach boxed title to top left={yshift=-2mm,xshift=4mm},
  boxed title style={colback=accent},drop shadow,#1}
\newtcolorbox{namebox}[1][]{enhanced,colback=deep!5,colframe=deep,
  fonttitle=\bfseries\color{white},coltitle=white,title={What This Is Called},
  attach boxed title to top left={yshift=-2mm,xshift=4mm},
  boxed title style={colback=deep},#1}
\newtcolorbox{reflectionbox}[1][]{enhanced,breakable,colback=purple!5,colframe=purple,
  fonttitle=\bfseries\color{white},coltitle=white,title={Reflection},#1}
\newtcolorbox{historybox}[1][]{enhanced,breakable,colback=gold!8,
  colframe=gold!70!brown,fonttitle=\bfseries\color{deep},title={Historical Note},#1}
\newtcolorbox{deepquestion}[1][]{enhanced,colback=accent!5,colframe=accent!60,
  fonttitle=\bfseries\color{accent},title={A Question Worth Sitting With},#1}
\newtcolorbox{realworldbox}[1][]{enhanced,breakable,colback=green!5,
  colframe=green!70!black,fonttitle=\bfseries\color{white},coltitle=white,
  title={In the Real World},#1}

\newcounter{problem}[chapter]
\newcommand{\problem}[1][]{\refstepcounter{problem}\medskip\noindent
  \textcolor{accent}{\textbf{Problem \theproblem.}}\ \ifx&#1&\else\textit{(#1)}\ \fi}
\newcommand{\keyidea}[1]{\begin{tcolorbox}[enhanced,colback=teal!8,colframe=teal,
  left=4mm,fontupper=\bfseries\color{deep}]#1\end{tcolorbox}}

\setlength{\epigraphwidth}{0.65\textwidth}
\renewcommand{\epigraphflush}{center}
\setlength{\parskip}{6pt}
\setlength{\parindent}{0pt}
\onehalfspacing

\begin{document}

\begin{titlepage}
\pagecolor{deep}\color{white}\centering\vspace*{1.8cm}

\begin{tikzpicture}[scale=1.2]
  % Symmetry and Noether's Theorem representation
  % Continuous translation of a wave
  \draw[teal!40, dashed, line width=1.2pt, domain=0:5, samples=100] plot (\x, {0.6*sin(360*\x/2.5) + 1.5});
  \draw[accent, line width=2pt, domain=0:5, samples=100] plot (\x, {0.6*sin(360*(\x-0.8)/2.5) + 1.5});
  
  \draw[->, gold, line width=2pt] (1.5, 2.3) -- ++(0.8, 0);
  \node[gold, font=\small] at (1.9, 2.6) {Translation $\Delta x$};
  
  \node[teal] at (2.5, 0.4) {Symmetry $\implies$ Conservation Law};
  \node[deep!30] at (2.5, -0.4) {Noether's Theorem};
\end{tikzpicture}

\vspace{0.8cm}
{\fontsize{11}{14}\selectfont\textcolor{teal}{\MakeUppercase{Discovery Physics · Journey 6}}}
\vspace{0.3cm}
{\fontsize{34}{40}\selectfont\bfseries Module 6.4}
\vspace{0.3cm}
{\fontsize{20}{26}\selectfont\textcolor{gold}{%
Quantum Operators and\\
Symmetries}}
\vspace{0.5cm}
\textcolor{white!60}{\rule{8cm}{0.4pt}}
\vspace{0.4cm}
{\fontsize{12}{18}\selectfont\textcolor{white!80}{\itshape
The algebra of observables.\\[4pt]
How do operators, commutators, and symmetries\\
dictate the conservation laws of nature?}}
\vfill\textcolor{white!30}{\small Version 1.0}
\end{titlepage}
\nopagecolor
\tableofcontents

\chapter{Quantum Operators and Symmetries}
\markboth{Journey 6: Module 6.4}{}

\begin{epigraph}
{Every continuous symmetry of a physical system corresponds to a conservation law of the system, and vice versa.}
{--- Emmy Noether}
\end{epigraph}

\section{Part 1 --- The Puzzle}

\begin{puzzlebox}
\textbf{Position vs. Momentum.}

In classical mechanics, position $x$ and momentum $p$ are simple real numbers. We can measure both simultaneously to infinite precision. In quantum mechanics, position and momentum are represented by linear operators: $\hat{x} = x$ and $\hat{p} = -i\hbar\frac{\partial}{\partial x}$.
* If we apply these operators to a wave function $\psi(x)$ in different orders, we get different results: measuring position then momentum is not the same as measuring momentum then position.
* This non-commutativity ($\hat{x}\hat{p} - \hat{p}\hat{x} \neq 0$) leads directly to Heisenberg's Uncertainty Principle:
  \[
    \sigma_x \sigma_p \;\ge\; \frac{\hbar}{2}.
  \]

\medskip
\textbf{Question 1.} Is this uncertainty a limitation of our experimental equipment? Or is it a fundamental mathematical property of wave physics, completely independent of measurement?

\textbf{Question 2.} In classical mechanics, we say a system has "symmetry under translation in space" if moving the whole system to the left doesn't change its behavior. Classically, this symmetry implies that total momentum is conserved. How does this link between symmetry and conservation manifest in the algebra of quantum operators?

\textbf{Question 3.} (Noether's Paradox.) Emmy Noether proved that *every* continuous symmetry corresponds to a conservation law. What quantum operator generates spatial translations? What operator generates time translations? How does the algebraic structure of these operators dictate the physics of the universe?
\end{puzzlebox}

\newpage

% ============================================================
\section{Part 2 --- Exploration}
% ============================================================

\subsection{Step 1: Operators and Expectation Values}

Every physical observable corresponds to a linear operator acting on the wave function.

\begin{explobox}
\problem Let $A$ be a physical property (like position, momentum, or energy). It is represented by a linear operator $\hat{A}$.
\begin{enumerate}[label=(\alph*), topsep=2pt]
  \item The expectation value $\langle A \rangle$ of the property in a state $\psi(x)$ is:
  \[
    \langle A \rangle \;=\; \int_{-\infty}^\infty \psi^* \hat{A} \psi \, dx.
  \]
  Because physical measurements must be real numbers, we must have $\langle A \rangle = \langle A \rangle^*$ for all states. Show that this requires the operator to be **Hermitian**:
  \[
    \int_{-\infty}^\infty \psi_1^* (\hat{A} \psi_2) \, dx \;=\; \int_{-\infty}^\infty (\hat{A} \psi_1)^* \psi_2 \, dx.
  \]
  \item Hermitian operators have two vital properties (linking back to Math Journey 3: Linear Structure): (i) their eigenvalues are strictly real, and (ii) eigenvectors corresponding to distinct eigenvalues are orthogonal. Explain why these properties are physically necessary for quantum measurements.
\end{enumerate}
\end{explobox}

\subsection{Step 2: Commutators and the Canonical Commutation Relation}

We study the algebraic structure of position and momentum operators.

\begin{explobox}
\problem The **commutator** of two operators $\hat{A}$ and $\hat{B}$ is defined as $[\hat{A}, \hat{B}] = \hat{A}\hat{B} - \hat{B}\hat{A}$.
\begin{enumerate}[label=(\alph*), topsep=2pt]
  \item Let $\hat{x} = x$ and $\hat{p} = -i\hbar\frac{d}{dx}$. We evaluate the commutator $[\hat{x}, \hat{p}]$ by letting it act on a test function $\psi(x)$:
  \[
    [\hat{x}, \hat{p}]\psi(x) \;=\; \hat{x}(\hat{p}\psi) - \hat{p}(\hat{x}\psi).
  \]
  Expand this expression using the product rule for the derivative term $\frac{d}{dx}(x\psi)$.
  \item Show that:
  \[
    [\hat{x}, \hat{p}]\psi(x) \;=\; i\hbar \psi(x).
  \]
  Conclude the **Canonical Commutation Relation**: $[\hat{x}, \hat{p}] = i\hbar$.
  \item If two operators commute ($[\hat{A}, \hat{B}] = 0$), they share a common basis of eigenvectors, meaning both properties can be measured simultaneously to infinite precision. Why does the fact that $[\hat{x}, \hat{p}] = i\hbar \neq 0$ guarantee that position and momentum can never be measured simultaneously?
\end{enumerate}
\end{explobox}

\subsection{Step 3: Symmetries and Generators}

We discover how quantum operators generate continuous physical translations.

\begin{explobox}
\problem Let us define a translation operator $\hat{T}(\Delta x)$ that shifts a wave function in space: $\hat{T}(\Delta x)\psi(x) = \psi(x - \Delta x)$.
\begin{enumerate}[label=(\alph*), topsep=2pt]
  \item For an infinitesimal translation $\Delta x \to 0$, Taylor-expand $\psi(x - \Delta x)$ to first order around $x$:
  \[
    \psi(x - \Delta x) \;\approx\; \psi(x) - \Delta x \frac{d\psi}{dx}.
  \]
  \item Substitute the momentum operator $\hat{p} = -i\hbar \frac{d}{dx}$ into the derivative term to show that:
  \[
    \hat{T}(\Delta x)\psi(x) \;\approx\; \left( 1 - \frac{i \Delta x}{\hbar} \hat{p} \right) \psi(x).
  \]
  \item We call $\hat{p}$ the **generator** of spatial translations.
  \item By the same logic, what operator generates translation in time ($\psi(t - \Delta t)$)? \textit{Hint: Look at the time derivative in the TDSE: $i\hbar \frac{\partial \Psi}{\partial t} = \hat{H}\Psi$}.
\end{enumerate}
\end{explobox}

\subsection{Step 4: Ehrenfest's Theorem and Conservation Laws}

We discover the link between operator commutation and physical conservation.

\begin{explobox}
\problem We study how expectation values change over time.
\begin{enumerate}[label=(\alph*), topsep=2pt]
  \item **Ehrenfest's Theorem** states that the time derivative of any expectation value is governed by its commutator with the Hamiltonian:
  \[
    \frac{d\langle A \rangle}{dt} \;=\; \frac{i}{\hbar} \langle [\hat{H}, \hat{A}] \rangle + \left\langle \frac{\partial \hat{A}}{\partial t} \right\rangle.
  \]
  \item If an operator has no explicit time dependence ($\partial \hat{A}/\partial t = 0$), show that $\langle A \rangle$ is constant in time (conserved) if and only if it commutes with the Hamiltonian:
  \[
    [\hat{H}, \hat{A}] \;=\; 0.
  \]
  \item **Noether's Theorem Connection:** If a system is invariant under spatial translations, its Hamiltonian $\hat{H}$ must commute with the generator of translations, which is the momentum operator $\hat{p}$.
  Show that if $[\hat{H}, \hat{p}] = 0$, then momentum is conserved: $\frac{d\langle p \rangle}{dt} = 0$.
  Symmetry under translation in space *is* the conservation of momentum.
\end{enumerate}
\end{explobox}

\newpage

% ============================================================
\section{Part 3 --- Discovery}
% ============================================================

\begin{discoverybox}
\textbf{What have we built?}

\textbf{First.} Physical observables are represented by Hermitian operators. Their eigenvalues are real (representing possible measurements) and their eigenvectors are orthogonal.

\textbf{Second.} The position and momentum operators do not commute: $[\hat{x}, \hat{p}] = i\hbar$. This canonical commutation relation dictates that measuring one observable disrupts the wave function's configuration in the other's basis.

\textbf{Third.} Momentum is the algebraic generator of spatial translations: $\hat{T}(\Delta x) \approx 1 - \frac{i\Delta x}{\hbar}\hat{p}$. The Hamiltonian (energy operator) is the generator of time translations.

\textbf{Fourth.} Under Ehrenfest's Theorem, an observable's expectation value is conserved if its operator commutes with the Hamiltonian. Noether's Theorem emerges in quantum mechanics as: \emph{A physical symmetry exists if the Hamiltonian commutes with the symmetry's generator, which guarantees the conservation of that generator.}
\end{discoverybox}

\newpage

% ============================================================
\section{Part 4 --- What This Is Called}
% ============================================================

\begin{namebox}
\textbf{Quantum Operators and Symmetries.}

* **Hermitian Operator:** An operator satisfying $\langle \psi_1 \mid \hat{A} \psi_2 \rangle = \langle \hat{A} \psi_1 \mid \psi_2 \rangle$.
* **Commutator:** $[\hat{A}, \hat{B}] = \hat{A}\hat{B} - \hat{B}\hat{A}$.
* **Canonical Commutation Relation:** $[\hat{x}, \hat{p}] = i\hbar$.
* **Ehrenfest's Theorem:**
  \[
    \frac{d\langle A \rangle}{dt} \;=\; \frac{i}{\hbar} \langle [\hat{H}, \hat{A}] \rangle + \left\langle \frac{\partial \hat{A}}{\partial t} \right\rangle.
  \]
* **Generator of Translation:** Momentum $\hat{p}$ generates translation in space; the Hamiltonian $\hat{H}$ generates translation in time.
* **Noether's Theorem (Quantum):** If $[\hat{H}, \hat{G}] = 0$, the physical property $G$ is conserved, and the system is symmetric under the transformations generated by $\hat{G}$.
\end{namebox}

\newpage

% ============================================================
\section{Part 5 --- Solved Exercises}
% ============================================================

\begin{exercise}[Hermiticity of the Momentum Operator]
Prove that the momentum operator $\hat{p} = -i\hbar \frac{d}{dx}$ is a Hermitian operator. You may assume the wave functions $\psi_1(x)$ and $\psi_2(x)$ vanish at infinity: $\lim_{x\to\pm\infty} \psi(x) = 0$.
\end{exercise}

\begin{proof}[Solution]
To prove Hermiticity, we must verify that:
\[
  \int_{-\infty}^\infty \psi_1^* (\hat{p} \psi_2) \, dx \;=\; \int_{-\infty}^\infty (\hat{p} \psi_1)^* \psi_2 \, dx.
\]
Let us write out the left-hand side, substituting $\hat{p} = -i\hbar \frac{d}{dx}$:
\[
  \text{LHS} \;=\; \int_{-\infty}^\infty \psi_1^* \left( -i\hbar \frac{d\psi_2}{dx} \right) \, dx \;=\; -i\hbar \int_{-\infty}^\infty \psi_1^* \frac{d\psi_2}{dx} \, dx.
\]
We use integration by parts: let $u = \psi_1^*$ and $dv = \frac{d\psi_2}{dx} dx$, so $du = \frac{d\psi_1^*}{dx} dx$ and $v = \psi_2$.
\[
  \text{LHS} \;=\; -i\hbar \left( \left[ \psi_1^* \psi_2 \right]_{-\infty}^\infty - \int_{-\infty}^\infty \psi_2 \frac{d\psi_1^*}{dx} \, dx \right).
\]
Since the wave functions vanish at infinity, the boundary term is 0: $\left[ \psi_1^* \psi_2 \right]_{-\infty}^\infty = 0 - 0 = 0$. Thus:
\[
  \text{LHS} \;=\; i\hbar \int_{-\infty}^\infty \psi_2 \frac{d\psi_1^*}{dx} \, dx \;=\; \int_{-\infty}^\infty \left( i\hbar \frac{d\psi_1}{dx} \right)^* \psi_2 \, dx.
\]
Observe that $i\hbar \frac{d\psi_1}{dx} = \left( -i\hbar \frac{d\psi_1}{dx} \right)^* = (\hat{p} \psi_1)^*$.
Substituting this back:
\[
  \text{LHS} \;=\; \int_{-\infty}^\infty (\hat{p} \psi_1)^* \psi_2 \, dx \;=\; \text{RHS}.
\]
This completes the proof. The momentum operator is Hermitian, guaranteeing that momentum measurements are always real numbers.
\end{proof}

\begin{exercise}[Ehrenfest's Theorem for Position]
Use Ehrenfest's Theorem to prove that the rate of change of expected position is related to expected momentum by:
\[
  \frac{d\langle x \rangle}{dt} \;=\; \frac{\langle p \rangle}{m}.
\]
(Note: The Hamiltonian is $\hat{H} = \frac{\hat{p}^2}{2m} + V(x)$).
\end{exercise}

\begin{proof}[Solution]
By Ehrenfest's Theorem, since the position operator $\hat{x} = x$ has no explicit time dependence ($\partial \hat{x}/\partial t = 0$):
\[
  \frac{d\langle x \rangle}{dt} \;=\; \frac{i}{\hbar} \langle [\hat{H}, \hat{x}] \rangle.
\]
Let us calculate the commutator $[\hat{H}, \hat{x}]$:
\[
  [\hat{H}, \hat{x}] \;=\; \left[ \frac{\hat{p}^2}{2m} + V(x), \, \hat{x} \right] \;=\; \frac{1}{2m} [\hat{p}^2, \hat{x}] + [V(x), \hat{x}].
\]
Since $V(x)$ and $\hat{x}$ are both simple multiplication operators in coordinate space, they commute: $[V(x), \hat{x}] = 0$.
Now calculate $[\hat{p}^2, \hat{x}]$ using the identity $[\hat{A}\hat{B}, \hat{C}] = \hat{A}[\hat{B}, \hat{C}] + [\hat{A}, \hat{C}]\hat{B}$:
\[
  [\hat{p}^2, \hat{x}] \;=\; \hat{p}[\hat{p}, \hat{x}] + [\hat{p}, \hat{x}]\hat{p}.
\]
Recall the canonical commutation relation: $[\hat{x}, \hat{p}] = i\hbar \implies [\hat{p}, \hat{x}] = -i\hbar$. Thus:
\[
  [\hat{p}^2, \hat{x}] \;=\; \hat{p}(-i\hbar) + (-i\hbar)\hat{p} \;=\; -2i\hbar \hat{p}.
\]
Substitute this back into the Hamiltonian commutator:
\[
  [\hat{H}, \hat{x}] \;=\; \frac{1}{2m} (-2i\hbar \hat{p}) \;=\; -\frac{i\hbar}{m} \hat{p}.
\]
Now substitute this commutator into Ehrenfest's Theorem:
\[
  \frac{d\langle x \rangle}{dt} \;=\; \frac{i}{\hbar} \left\langle -\frac{i\hbar}{m} \hat{p} \right\rangle \;=\; \frac{i}{\hbar} \left(-\frac{i\hbar}{m}\right) \langle p \rangle \;=\; \frac{\langle p \rangle}{m}.
\]
This completes the proof. This result shows that quantum expected values obey the classical relation $v = p/m$.
\end{proof}

\newpage

% ============================================================
\section{Part 6 --- Problems}
% ============================================================

\begin{problem}[Uncertainty as a Wave Property]
The wave function in momentum space $\phi(p)$ is the Fourier transform of the spatial wave function $\psi(x)$:
\[
  \phi(p) \;=\; \frac{1}{\sqrt{2\pi\hbar}} \int_{-\infty}^\infty \psi(x) e^{-i p x/\hbar} \, dx.
\]
\begin{enumerate}[label=(\alph*)]
  \item Explain why a highly localized particle in space (a narrow spike in $\psi(x)$) requires a wide range of plane-wave frequencies (a broad spread in $\phi(p)$) to reconstruct it.
  \item Why does this Fourier relationship mean that Heisenberg's Uncertainty Principle $\sigma_x \sigma_p \ge \hbar/2$ is a fundamental property of waves, completely independent of any physical measurement disturbance?
\end{enumerate}
\end{problem}

\begin{problem}[Angular Momentum Commutators]
In 3D, the orbital angular momentum operator around the $z$-axis is:
\[
  \hat{L}_z \;=\; \hat{x}\hat{p}_y - \hat{y}\hat{p}_x.
\]
\begin{enumerate}[label=(\alph*)]
  \item Prove the commutator relations:
  \[
    [\hat{L}_z, \hat{x}] \;=\; i\hbar \hat{y}, \qquad [\hat{L}_z, \hat{y}] \;=\; -i\hbar \hat{x}.
  \]
  \textit{Hint: Recall that operators in different coordinates commute (e.g., $[\hat{x}, \hat{p}_y] = 0$).}
  \item Explain why these commutators show that we cannot simultaneously measure angular momentum around the $z$-axis and coordinate positions in the $xy$-plane.
\end{enumerate}
\end{problem}

\newpage

% ============================================================
\section{Part 7 --- Deep Problems}
% ============================================================

\begin{deepproblem}[Time Translation and Energy Conservation]
We study the symmetry of a system under translation in time.
\begin{enumerate}[label=(\alph*)]
  \item Suppose a system's Hamiltonian does not depend explicitly on time: $\partial\hat{H}/\partial t = 0$.
  Use Ehrenfest's Theorem to prove that the expected energy is constant in time:
  \[
    \frac{d\langle H \rangle}{dt} \;=\; 0.
  \]
  \item Under Noether's Theorem, this means that invariance under translation in time *is* the conservation of energy. Explain this relationship in your own words.
\end{enumerate}
\end{deepproblem}

\begin{deepproblem}[The Generalized Uncertainty Principle]
Let $\hat{A}$ and $\hat{B}$ be two Hermitian operators. The variance of measurements is $\sigma_A^2 = \langle (\hat{A} - \langle A \rangle)^2 \rangle$ and $\sigma_B^2 = \langle (\hat{B} - \langle B \rangle)^2 \rangle$.
Prove the **Robertson Uncertainty Relation**:
\[
  \sigma_A^2 \sigma_B^2 \;\ge\; \left| \frac{1}{2i} \langle [\hat{A}, \hat{B}] \rangle \right|^2.
\]
\textit{Hint: Let $\hat{A}' = \hat{A} - \langle A \rangle$ and $\hat{B}' = \hat{B} - \langle B \rangle$. Use the Cauchy-Schwarz inequality for integrals $\langle f \mid f \rangle \langle g \mid g \rangle \ge |\langle f \mid g \rangle|^2$ with $f = \hat{A}'\psi$ and $g = \hat{B}'\psi$, and split the inner product $\langle \hat{A}'\psi \mid \hat{B}'\psi \rangle$ into its symmetric and anti-symmetric parts.}
\end{deepproblem}

\newpage

% ============================================================
\section{Part 8 --- Reflection}
% ============================================================

\begin{reflectionbox}
\textbf{Conceptual Reflection}

Consider the following questions:
\begin{enumerate}
  \item **On Symmetries:** Noether's theorem shows that symmetries of space and time dictate the conservation laws of our universe (translation in space $\implies$ momentum; translation in time $\implies$ energy; rotation $\implies$ angular momentum). Does this mean physics is ultimately a study of geometry?
  \item **On Heisenberg:** If the uncertainty principle is a mathematical property of wave packets, why did Heisenberg explain it originally as a physical disturbance caused by a microscope shining light on an electron?
  \item **On Operators:** In classical mechanics, observables are properties that a particle *possesses*. In quantum mechanics, observables are operators that *act* on the state. What does this transition suggest about the nature of physical properties in the quantum world?
\end{enumerate}
\end{reflectionbox}

\medskip
\noindent\textbf{Looking back at Journey 6.}
You have traversed the foundations of statistical mechanics and quantum theory:
* **Module 6.1:** Microstates, partition functions, and entropy emergence.
* **Module 6.2:** Blackbody mode densities, energy quantization, and Planck's Law.
* **Module 6.3:** de Broglie waves, the Schrödinger equation, and standing-wave bound energy levels.
* **Module 6.4:** Observable operators, commutators, symmetries, and Noether's conservation laws.

This completes Journey 6, providing the absolute physical frontier of the curriculum. You are now equipped with the mathematical and physical tools to understand solid-state band structures, chemical kinetics, and quantum information in the journeys ahead.

\vspace{1.4cm}
\begin{center}
{\color{teal}\rule{0.5\linewidth}{0.8pt}}\\[6pt]
{\small\color{deep}\textit{End of Journey 6 --- Statistical Mechanics & The Quantum Frontier}}\\[4pt]
{\small\color{deep}You have completed the Physics Frontier Track. Continue to: \textbf{Chemistry Journey 5 --- Chemical Kinetics}}\\[6pt]
{\color{teal}\rule{0.5\linewidth}{0.8pt}}
\end{center}

\end{document}
